\chapter{The Seiberg-Witten Equations and Moduli Space}

\dropcap{W}{e are finally} ready to define the Seiberg-Witten equations and their moduli space. Let $M$ be a closed oriented Riemannian four-manifold with a fixed $\Spinc$ structure. The equations are written for a spinor and a compatible connection.

\section{The gauge group and its action}
\begin{definitionbox}[Gauge group]
For a principal $U(1)$-bundle $P\to M$, the gauge group is
\[
  \cG := \operatorname{Aut}_{U(1)}(P).
\]
In many settings this may be identified with $C^\infty(M,U(1))$.
\end{definitionbox}

\section{Topology of the Moduli Space}
\begin{definitionbox}[Seiberg-Witten equations]
The Seiberg-Witten equations for a spinor $\psi$ and a connection $A$ are
\[
  \Dslash_A\psi=0,
  \qquad
  F_A^+ = \sigma^+(\psi).
\]
\end{definitionbox}

\begin{definitionbox}[Moduli space]
The Seiberg-Witten moduli space is the quotient
\[
\cM = \{(\psi,A)\mid \Dslash_A\psi=0,\;F_A^+=\sigma^+(\psi)\}/\cG.
\]
\end{definitionbox}

\subsection{Dimension of $\cM$}
The expected dimension is computed from an elliptic deformation complex and an index theorem.

\subsection{Generic smoothness}
For generic perturbations, the moduli space is smooth under suitable hypotheses.

\subsection{Orientation}
Orienting the moduli space requires orienting a determinant line associated to the deformation complex.

\subsection{Compactness}
Compactness follows from a priori estimates and elliptic regularity.

\section{The Seiberg-Witten invariant}
When the moduli space has the correct dimension, one can extract invariants of the smooth structure of the underlying four-manifold.
